A predição e classificação de POIs são desafios fundamentais em sistemas de recomendação baseados em localização e em análise de mobilidade urbana. Esta seção organiza os trabalhos relacionados em quatro eixos: representações baseadas em grafos, codificadores espaciais, representações temporais e aprendizado multitarefa aplicado a POIs.

\subsection{Representações de POIs Baseadas em Grafos}

A modelagem de POIs por meio de grafos permite capturar relações de vizinhança espacial e semântica entre locais. O modelo POI2Vec \cite{feng2017poi2vec} adapta a arquitetura Word2Vec para gerar \textit{embeddings} de POIs que incorporam a influência geográfica entre locais por meio de uma estrutura de árvore binária geográfica. O CATAPE (\textit{Category-Aware Location Embedding}) \cite{rahmani2019category} estende essa ideia incorporando informações categóricas e sequenciais, capturando a influência geográfica entre POIs com base na sequência temporal de visitas dos usuários.

Em uma abordagem auto-supervisionada, o \textit{Deep Graph Infomax} (DGI) \cite{velickovic2019deep} maximiza a informação mútua entre representações locais (nós) e globais (grafo) para gerar \textit{embeddings} sem supervisão explícita. O HGI (\textit{Hierarchical Graph Infomax}) \cite{huang2023learning} avança sobre essa linha ao operar em múltiplos níveis hierárquicos (POI, região e cidade), produzindo representações enriquecidas com informações regionais por meio de convoluções em grafos e mecanismos de multi-atenção. 
% O HAVANA \cite{dos2024havana} combina \textit{Graph Attention Networks} (GAT) \cite{velivckovic2017graph} com filtros ARMA para anotação semântica de locais, mostrando a eficácia de modelos híbridos de grafos para a tarefa de classificação categórica.

No entanto, essas abordagens codificam informações espaciais de forma implícita por meio da topologia do grafo, sem utilizar diretamente as coordenadas geográficas como sinal de entrada para a representação.

\subsection{Representações Espaciais}

O uso direto de coordenadas geográficas para a geração de características latentes espaciais é uma abordagem distinta das representações baseadas em grafos. O \textit{framework} TorchSpatial \cite{wu2024torchspatial} estabelece um \textit{benchmark} unificado para a avaliação de estratégias de representação espacial, permitindo a comparação sistemática entre diferentes arquiteturas.

% O Space2Vec \cite{mai2020multiscalerepresentationlearningspatial} propõe uma codificação posicional multi-escala inspirada em \textit{grid cells} do sistema nervoso, utilizando funções senoidais com frequências logaritmicamente espaçadas para transformar coordenadas em representações vetoriais. A abordagem 
A metodologia SIREN (\textit{Sinusoidal Representation Networks}) \cite{sitzmann2020implicit}, aplicada ao contexto geográfico \cite{rußwurm2024geographiclocationencodingspherical}, modela funções contínuas por meio de ativações senoidais, sendo capaz de representar sinais de alta frequência em dados espaciais. O Sphere2Vec \cite{mai2023sphere2vecgeneralpurposelocationrepresentation} opera diretamente sobre coordenadas esféricas com múltiplas escalas de resolução, preservando propriedades de distância geodésica que são perdidas em projeções planares.

Embora essas estratégias tenham sido avaliadas em tarefas geoespaciais como predição de espécies e estimativa de população, sua aplicação pode ser estendida para modelos de predição e classificação de POIs como proposto neste trabalho.

\subsection{Representações Temporais}

A dimensão temporal é central para tarefas de predição de POIs, uma vez que os padrões de visitação dos usuários apresentam regularidades (e.g., horários de refeição, deslocamentos semanais). O modelo LSTPM (\textit{Long- and Short-Term Preference Modeling}) \cite{sun2020go} captura preferências de curto e longo prazo utilizando uma Geo-Dilated RNN que relaciona POIs não consecutivos por meio de sequências temporais de visita e informações geográficas. O TransTARec \cite{sun2024transtarec} propõe um modelo de tradução temporal adaptativa que une representações do POI, do \textit{timestamp} e das preferências do usuário para predição do próximo POI.

Esses modelos incorporam a dimensão temporal principalmente no contexto da sequência de visitas do usuário. Nesse sentido, o Time2Vec \cite{kazemi2019time2vec} propõe uma representação temporal que combina termos lineares com funções senoidais de frequência e fase aprendíveis, capturando padrões periódicos e não periódicos a partir de valores temporais contínuos. 

\subsection{Aprendizado Multitarefa para Tarefas de POI}

O aprendizado multitarefa (\textit{Multitask Learning}, MTL) \cite{caruana1997multitask} tem sido explorado em tarefas de POIs principalmente para combinar a predição do próximo POI com sinais temporais, comportamentais ou regionais. O MCARNN \cite{Liao2018} integra dependências sequenciais e regularidades temporais para predizer simultaneamente atividades e locais futuros. O MTPR \cite{Xia2020} modela conjuntamente localização e contexto temporal por meio de LSTMs geográficas e aprendizado adversarial. Em uma linha mais recente, Halder et al. \cite{Halder2022} propõem um modelo multitarefa baseado em \textit{Transformers} para recomendar o próximo POI e predizer simultaneamente o tempo de espera.

Na classificação categórica, o TME \cite{Xu2023} explora uma estrutura hierárquica de categorias para anotação semântica de POIs, enquanto o HMT-GRN \cite{Lim2022} combina predição do próximo POI e de sua região geográfica em um modelo recorrente com grafos hierárquicos. Em conjunto, esses trabalhos indicam que o MTL é promissor em problemas de POIs, mas ainda há pouca exploração de arquiteturas que tratem conjuntamente a \textit{classificação de categoria de POI} e a \textit{predição do próximo POI} com representações espaciais contínuas e temporais explícitas. Essa lacuna motiva a investigação proposta neste trabalho.

